% !TEX root = task-3-5-2.tex
% Editable task bank. Append new BankTask blocks at the end of the document.
\documentclass[12pt,a4paper]{article}
\usepackage{physics-textbook}
\usepackage{graphicx,needspace}
\usetikzlibrary{decorations.pathmorphing}
\newif\ifShowSolutions
\ifdefined\TutorEdition
  \ShowSolutionstrue
\else
  \ShowSolutionsfalse % Change to \ShowSolutionstrue to display solutions.
\fi
\newcounter{banktask}
\newsavebox{\BankTaskBox}
\newcommand{\BankTaskContent}[3]{%
  \noindent\textbf{\thebanktask.}\quad #2\par
  \ifShowSolutions
    \Needspace{50mm}% Keep the solution heading with its opening text or diagram.
    #3\par
    \ifdefined\TutorEdition
      \ifcsname TutorSource#1\endcsname
        \nopagebreak\smallskip{\small\csname TutorSource#1\endcsname\par}%
      \fi
    \fi
  \fi
}
\newcommand{\BankTask}[3]{%
  \refstepcounter{banktask}\par\medskip
  \sbox{\BankTaskBox}{\begin{minipage}{\linewidth}%
    \BankTaskContent{#1}{#2}{#3}\end{minipage}}%
  \ifdim\dimexpr\ht\BankTaskBox+\dp\BankTaskBox\relax>\dimexpr\textheight-20mm\relax
    \Needspace{35mm}\BankTaskContent{#1}{#2}{#3}%
  \else
    \Needspace{\dimexpr\ht\BankTaskBox+\dp\BankTaskBox\relax}%
    \noindent\usebox{\BankTaskBox}\par
  \fi
  \medskip
}
\begin{document}
\BookSubsection{3.5.2}{Заряд внутри проводника}


% Source sdamgia-9088: https://phys-ege.sdamgia.ru/problem?id=9088
\BankTask{sdamgia-9088}{
Если к незаряженному металлическому шару поднести, не касаясь, точечный положительный заряд, то на стороне шара, ближайшей к заряду, появится отрицательный заряд. Как называется это явление (электризация, электростатическая индукция, электромагнитная индукция, поляризация)?\par
}{
\smallskip\noindent\textbf{Решение}\par
Явление наведения собственного электростатического поля при действии на тело внешнего электрического поля называется электростатической индукцией.\par
Ответ: электростатическая индукция.\par
}

% Source sdamgia-31415: https://phys-ege.sdamgia.ru/problem?id=31415
\BankTask{sdamgia-31415}{
Стеклянные незаряженные кубики~$1$ и $2$ и алюминиевые незаряженные кубики~$3$ и $4$ сблизили вплотную и поместили в электрическое поле, напряженность которого направлена горизонтально влево, как показано в левой части рисунка. Затем кубики раздвинули и уже после этого выключили электрическое поле (правая часть рисунка).\par
$1$. В электрическом поле (левая часть рисунка) часть электронов перешла из кубика~$4$ в кубик~$3$.\par
$2$. Кубик $2$, изображенный в правой части рисунка, не заряжен.\par
$3$. В электрическом поле (левая часть рисунка) правая грань кубика~$2$ заряжена положительно.\par
$4$. Кубик $3$, изображенный в правой части рисунка, имеет положительный заряд.\par
$5$. В электрическом поле (левая часть рисунка) разность потенциалов между левой гранью кубика~$3$ и правой гранью кубика~$4$ равна нулю.\par

\begin{center}\begin{tikzpicture}[book diagram,>=stealth]
\draw(0,0)rectangle(1.2,1.2);\draw(.6,0)--(.6,1.2)(0,.6)--(1.2,.6);\foreach\x/\y/\s in{.3/.9/1,.9/.9/2,.3/.3/3,.9/.3/4}\node at(\x,\y){$\s$};\foreach\y in{.3,.5,.7,.9}\draw[->](-.2,\y)--(-1,\y);\node[above]at(-.7,1){$\vec E$};\draw[dashed](1.65,-.2)--(1.65,1.4)(1.8,-.2)--(1.8,1.4);\foreach\x/\y/\s in{2.1/.8/1,3.1/.8/2,2.1/-.2/3,3.1/-.2/4}{\draw(\x,\y)rectangle++(.6,.6);\node at(\x+.3,\y+.3){$\s$};}
\end{tikzpicture}\end{center}
\par
}{
\smallskip\noindent\textbf{Решение}\par
$1$. Неверно. Электроны под действием электрического поля движутся против силовых линий, поэтому переместятся вправо, тогда часть электронов перейдет с кубика~$3$ в кубик~$4$.\par
$2$. Верно. Стекло является диэлектриком, в котором нет свободных носителей заряда. Поэтому кубики из стекла останутся незаряженными.\par
$3$. Неверно. В кубиках из стекла произойдет поляризация, в результате которой правая часть кубика~$2$ будет иметь отрицательный заряд.\par
$4$. Верно. Поскольку электроны в металлах переместятся против силовых линий, то в кубике~$3$ будет недостаток электронов, он будет заряжен положительно.\par
$5$. Верно. В проводнике, помещенном в электростатическое поле, напряженность станет равной $0$, поэтому и разность потенциалов между левой гранью кубика~$3$ и правой гранью кубика~$4$ равна нулю.\par
Ответ: $245$.\par
}

% Source sdamgia-6431: https://phys-ege.sdamgia.ru/problem?id=6431
\BankTask{sdamgia-6431}{
В нижней половине незаряженного металлического шара находится крупная шарообразная полость, заполненная воздухом. Шар находится в воздухе вдали от других предметов. В центр полости помещен положительный точечный заряд $q > 0$ (см. рис.). Нарисуйте картину силовых линий электростатического поля внутри полости и снаружи шара. Если поле равно нулю, напишите в данной области: $\vec{E}=0.$ Если поле отлично от нуля, нарисуйте картину поля в данной области, используя восемь силовых линий.\par

\begin{center}\begin{tikzpicture}[book diagram,>=stealth]
\draw[fill=gray!12](0,0)circle(1.5);\draw[fill=white](0,-.5)circle(.75);\fill(0,-.5)circle(1.5pt);\node[below right]at(0,-.5){$q$};
\end{tikzpicture}\end{center}
\par
}{
\smallskip\noindent\textbf{Решение}\par

\begin{center}\begin{tikzpicture}[book diagram,>=stealth]
\draw[fill=gray!12](0,0)circle(1.5);\draw[fill=white](0,-.5)circle(.75);\node at(0,.8){$E=0$};\foreach\a in{0,45,...,315}{\draw[->](\a:1.5)--(\a:2.3);\draw[->](0,-.5)--++(\a:.72);}
\end{tikzpicture}\end{center}
\par
$1$. Приведен схематический рисунок картины силовых линий: внутри полости — семейство прямых лучей, исходящих из заряда и приходящих на поверхность полости по нормали; снаружи шара — семейство прямых лучей, исходящих с поверхности шара по нормали к ней и уходящих на бесконечность.\par
$2$. Внутри проводника — электростатическое поле $E=0.$ Поэтому поле в полости обладает центральной симметрией и выглядит как поле точечного заряда $q > 0,$ находящегося в центре полости. Силовые линии этого поля подходят по нормали к поверхности полости, где равномерно распределен отрицательный индуцированный заряд $-q < 0.$\par
$3$. На наружной поверхности шара находится (в силу нейтральности шара в целом) положительный заряд $q > 0.$ В силу того, что внутри проводника $\vec{E}=0,$ а снаружи окружающие предметы расположены далеко от шара, этот заряд распределен по поверхности шара равномерно. Его поле вне шара выглядит как поле точечного заряда $q > 0,$ расположенного в центре шара. Силовые линии отходят от шара по нормали к его поверхности.\par
}

% Source sdamgia-4218: https://phys-ege.sdamgia.ru/problem?id=4218
\BankTask{sdamgia-4218}{
Внутри незаряженного металлического шара радиусом $r_{1}=40 \,\text{см}$ имеются две сферические полости радиусами $r_{2} < \frac{ r_{1}}{ 2 } ,$ расположенные таким образом, что их поверхности почти соприкасаются в центре шара. В центре одной полости поместили заряд $q_{1}= + 1\,\text{нКл}$, а затем в центре другой — заряд $q_{2}= + 2\,\text{нКл}$ (см. рис.). Найдите модуль и направление вектора напряженности $\vec{E}$ электростатического поля в точке $O$, находящейся на расстоянии $R=1 \,\text{м}$ от центра шара на перпендикуляре к отрезку, соединяющему центры полостей.\par

\begin{center}\begin{tikzpicture}[book diagram,>=stealth]
\draw[fill=gray!12](0,0)circle(1.5);\draw[fill=white](0,.67)circle(.67)(0,-.67)circle(.67);\draw(-1.5,0)--(2.7,0);\draw(0,.67)--(0,-.67);\draw(0,0)--(.14,0)--(.14,.14)--(0,.14);\foreach\y/\s in{.67/q_1,-.67/q_2}{\fill(0,\y)circle(1.5pt);\node[right]at(0,\y){$\s$};}\fill(2.7,0)circle(1.5pt)node[above right]{$O$};\node[above]at(-.85,0){$r_1$};\node[left]at(0,.35){$r_2$};\node[left]at(0,-.35){$r_2$};\node[above]at(1.9,0){$R$};
\end{tikzpicture}\end{center}
\par
}{
\smallskip\noindent\textbf{Решение}\par
Возможное решение\par
В электростатике считается, что электрическое поле внутри металла отсутствует, так как иначе свободные заряды внутри металла двигались бы. Поэтому при помещении заряда $q_{1}$ в первую полость на ее стенках индуцируется заряд $- q_{1}$ и по принципу суперпозиции суммарное поле этих двух зарядов в металле шара равно нулю. По закону сохранения электрического заряда и в силу электронейтральности шара избыточный заряд равен $q_{1}.$ Он вытесняется на поверхность шара и равномерно распределяется по ней, так как заряды внутри металла не создают поля и не влияют на распределение зарядов на поверхности шара.\par
После помещения заряда $q_{2}$ во вторую полость все происходит аналогичным образом, и на поверхность шара вытесняется дополнительно заряд $q_{2},$ так что теперь на поверхности шара равномерно распределяется суммарный заряд $q_{1} + q_{2}.$\par
Поле равномерно заряженного шара вне его совпадает с полем точечного заряда, помещенного в центр шара, поэтому поле в точке $O$, находящейся на расстоянии $R$ от центра шара, по модулю равно $E= \frac{ 1}{ 4 \pi \varepsilon _{0} } \cdot \frac{ q_{1} + q_{2}}{ R ^2 } =27\,\text{В/м}$ и направлено вдоль радиуса от центра шара к точке $O.$\par
Здесь $\varepsilon _{0}$ — электрическая постоянная.\par
Ответ: $E=27\,\text{В/м}$, вектор $\vec{E}$ направлен от центра шара к точке $O.$\par
}

% Source sdamgia-6067: https://phys-ege.sdamgia.ru/problem?id=6067
\BankTask{sdamgia-6067}{
Если потереть шерстью эбонитовую палочку, то она электризуется, приобретая отрицательный заряд, и стрелка электрометра при поднесении палки к его шару отклоняется, а при удалении палки — возвращается к неотклоненному состоянию. Если же в момент поднесения наэлектризованной палки к электрометру коснуться рукой его металлического корпуса и сразу же убрать руку, то после удаления палки отклонение стрелки сохраняется, хотя и меньшее по величине.\par
Объясните, основываясь на известных физических законах и закономерностях, почему это происходит.\par
Электрометр (см. рис.) представляет собой металлический цилиндрический корпус, передняя и задняя стенки которого стеклянные. Корпус закреплен на изолирующей подставке. Через изолирующую втулку внутрь корпуса сверху входит металлическая трубка, заканчивающаяся внизу стержнем с установленной на нем легкоподвижной стрелкой, отклонение которой определяется величиной заряда. Стрелка может вращаться вокруг горизонтальной оси. Внутри корпуса установлена шкала электрометра, по которой определяется отклонение стрелки. Снаружи корпуса, наверху трубки прикрепляется металлический шар или тарелка, к которой подносят заряженные тела.\par

\begin{center}\begin{tikzpicture}[book diagram,>=stealth]
\begin{scope}[shift={(0,0)}]\draw(0,0)circle(.65);\draw(0,.65)--(0,1.1);\draw(0,1.3)circle(.2);\draw(0,-.65)--(0,-1)--(-.4,-1)(0,-1)--(.4,-1);\draw(0,.35)--(0,-.35);\draw[thick](0.4,.48)--(-0.4,-.48);\foreach\a in{220,235,...,280}\draw({.48*cos(\a)},{.48*sin(\a)})--({.58*cos(\a)},{.58*sin(\a)});\end{scope}\draw[thick](.4,1.7)--(1.7,1.2);
\end{tikzpicture}\end{center}
\par
}{
\smallskip\noindent\textbf{Решение}\par
$1$. При поднесении наэлектризованной отрицательным зарядом эбонитовой палки к шару электрометра, в силу явления электростатической индукции и закона сохранения заряда в изолированной системе тел, шар и стрелка электрометра заряжаются разноименными и равными по величине зарядами (шар — «+», стрелка — «-»). При этом часть металлического корпуса электрометра вблизи шкалы заряжается положительным зарядом в силу того же явления электростатической индукции, а остальная часть — равным ему по величине в силу закона сохранения заряда в изолированной системе тел отрицательным зарядом.\par
$2$. Стрелка электрометра отклоняется, так как одноименные заряды на стрелке и на стержне отталкиваются, а разноименные заряды на стрелке и на корпусе электрометра притягиваются, согласно закону взаимодействия зарядов.\par
$3$. При удалении наэлектризованной палки от электрометра одинаковые индуцированные заряды разных знаков на его шаре и на стрелке, а также меньшие по величине заряды на корпусе электрометра вблизи его шкалы и вдали от стрелки компенсируются, и отклонение стрелки прекращается.\par
$4$. Если коснуться корпуса электрометра рукой после поднесения к нему наэлектризованной палки и сразу убрать руку, то часть индуцированных на корпусе зарядов (отрицательных) стечет на прикоснувшегося к нему\par
человека, и на корпусе электрометра останется нескомпенсированный положительный заряд.\par
$5$. В силу явления электростатической индукции после удаления палки этот положительный нескомпенсированный заряд на корпусе электрометра вызовет появление отрицательного заряда на стрелке электрометра, расположенной вблизи шкалы, и положительного — на шаре электрометра, что и приведет к отклонению стрелки, хотя и меньшему, чем при поднесении заряженной эбонитовой палки к электрометру.\par
}

% Source sdamgia-38948: https://phys-ege.sdamgia.ru/problem?id=38948
\BankTask{sdamgia-38948}{
На столе установили два незаряженных электрометра и соединили их металлическим стержнем с изолирующей ручкой (рис. $1$). Затем к первому электрометру поднесли, не касаясь его шара, положительно заряженную палочку (рис. $2$). Не убирая палочки, сняли стержень, а затем убрали палочку.\par
Ссылаясь на известные вам законы и явления, объясните, почему электрометры оказались заряженными, и определите знаки заряда каждого из электрометров после того, как палочку убрали.\par

\begin{center}\begin{tikzpicture}[book diagram,>=stealth]
\begin{scope}[shift={(0,0)}]\draw(0,0)circle(.65);\draw(0,.65)--(0,1.1);\draw(0,1.3)circle(.2);\draw(0,-.65)--(0,-1)--(-.4,-1)(0,-1)--(.4,-1);\draw(0,.35)--(0,-.35);\draw[thick](0,.48)--(0,-.48);\foreach\a in{220,235,...,280}\draw({.48*cos(\a)},{.48*sin(\a)})--({.58*cos(\a)},{.58*sin(\a)});\end{scope}\begin{scope}[shift={(2.4,0)}]\draw(0,0)circle(.65);\draw(0,.65)--(0,1.1);\draw(0,1.3)circle(.2);\draw(0,-.65)--(0,-1)--(-.4,-1)(0,-1)--(.4,-1);\draw(0,.35)--(0,-.35);\draw[thick](0,.48)--(0,-.48);\foreach\a in{220,235,...,280}\draw({.48*cos(\a)},{.48*sin(\a)})--({.58*cos(\a)},{.58*sin(\a)});\end{scope}\draw(0,1.5)--(2.4,1.5);\draw[fill=white](1.1,1.35)rectangle(1.25,1.85);\node[below]at(0,-1){$1$};\node[below]at(2.4,-1){$2$};\node[below]at(1.2,-1.5){Рис. 1};
\end{tikzpicture}\end{center}
\par

\begin{center}\begin{tikzpicture}[book diagram,>=stealth]
\begin{scope}[shift={(0,0)}]\draw(0,0)circle(.65);\draw(0,.65)--(0,1.1);\draw(0,1.3)circle(.2);\draw(0,-.65)--(0,-1)--(-.4,-1)(0,-1)--(.4,-1);\draw(0,.35)--(0,-.35);\draw[thick](.52,.3)--(-.52,-.3);\foreach\a in{220,235,...,280}\draw({.48*cos(\a)},{.48*sin(\a)})--({.58*cos(\a)},{.58*sin(\a)});\end{scope}\begin{scope}[shift={(2.4,0)}]\draw(0,0)circle(.65);\draw(0,.65)--(0,1.1);\draw(0,1.3)circle(.2);\draw(0,-.65)--(0,-1)--(-.4,-1)(0,-1)--(.4,-1);\draw(0,.35)--(0,-.35);\draw[thick](.52,.3)--(-.52,-.3);\foreach\a in{220,235,...,280}\draw({.48*cos(\a)},{.48*sin(\a)})--({.58*cos(\a)},{.58*sin(\a)});\end{scope}\draw(0,1.5)--(2.4,1.5);\draw[fill=white](1.1,1.35)rectangle(1.25,1.85);\draw[thick](-1.1,1.7)--(-.25,2.1);\node[above]at(-.7,2){$+$};\node[below]at(0,-1){$1$};\node[below]at(2.4,-1){$2$};\node[below]at(1.2,-1.5){Рис. 2};
\end{tikzpicture}\end{center}
\par
}{
\smallskip\noindent\textbf{Решение}\par
Вокруг положительно заряженной палочки существует электрическое поле. Под его действием свободные электроны в проводнике (металлической палочке, металлических шарах) придут в движение, причем сила электрического действия будет направлена к заряженной палочке. В результате на шаре электрометра $1$ возникнет избыток электронов, а на шаре электрометра $2$ возникнет недостаток электронов. Данный процесс будет происходить до тех пор, пока потенциалы всех точек не станут равными. Таким образом, электрометр $1$ будет заряжен отрицательно, а электрометр $2$ — положительно.\par
Электрометры образуют изолированную систему, в которой выполняется закон сохранения электрического заряда. Так как до поднесения палочки система не имела заряда, то после на каждом шаре заряды будут равными по модулю и противоположными по знаку. После удаления стержня и палочки на каждом шаре останется индуцированный заряд.\par
Ответ: электрометр $1$ имеет отрицательный заряд, электрометр $2$ — положительный.\par
}


% Source 896222: https://ege.fipi.ru/bank/questions.php?search=1&proj=BA1F39653304A5B041B656915DC36B38&qid=896222
\BankTask{896222}{
Незаряженное металлическое тело, продольное сечение которого показано на рисунке, поместили в однородное электрическое поле напряжённостью $\vec{E}$.

Какие два утверждения описывают результаты воздействия этого поля на металлическое тело?
\begin{enumerate}
\item Напряжённость электрического поля в точке $D$ равна нулю.
\item Потенциалы в точках $B$ и $C$ равны.
\item Концентрация свободных электронов в точке $A$ наибольшая.
\item В точке $A$ индуцируется положительный заряд.
\item В точке $D$ индуцируется отрицательный заряд.
\end{enumerate}
\begin{center}
\begin{tikzpicture}[book diagram,x=1cm,y=1cm]

\draw[-{book tip}] (0,2.45)--(3.8,2.45);
\draw[-{book tip}] (0,0.15)--(3.8,0.15);
\node[above] at (3.6,2.55) {$\vec E$};
\node[above] at (3.6,0.25) {$\vec E$};
\draw[fill=gray!55,line width=0.6pt] (0.85,0.87)--(3.05,0.87) arc[start angle=-90,end angle=90,radius=0.45]--(0.85,1.77) arc[start angle=90,end angle=270,radius=0.45]--cycle;
\fill (1.95,1.77) circle (0.035);
\fill (0.4,1.32) circle (0.035);
\fill (3.5,1.32) circle (0.035);
\fill (1.95,1.36) circle (0.04);
\node[above] at (1.95,1.82) {$A$};
\node[left] at (0.4,1.32) {$B$};
\node[right] at (3.5,1.32) {$C$};
\node[right] at (1.97,1.3) {$D$};
\end{tikzpicture}
\end{center}
}{
\smallskip\noindent\textbf{Решение от GPT}\par
В электростатическом равновесии внутри металла $\vec E=0$, а весь проводник эквипотенциален: $\varphi_B=\varphi_C$; утверждения 1 и 2 верны. Электроны смещаются против внешнего поля и накапливаются у левого края $B$, поэтому утверждение 3 неверно. В точке $A$ на середине боковой поверхности положительный заряд не индуцируется, а внутри металла в точке $D$ избыточного заряда нет, поэтому утверждения 4 и 5 неверны. Ответ: $12$.
}

% Source E412A6: https://ege.fipi.ru/bank/questions.php?search=1&proj=BA1F39653304A5B041B656915DC36B38&qid=E412A6
\BankTask{E412A6}{
Стеклянные незаряженные кубики 1 и 2 и медные незаряженные кубики 3 и 4 сблизили вплотную и поместили в электрическое поле, напряжённость которого направлена горизонтально вправо, как показано в левой части рисунка. Затем кубики раздвинули и уже после этого выключили электрическое поле (правая часть рисунка).

Какие два утверждения описывают данный процесс?
\begin{enumerate}
\item В электрическом поле кубики 1 и 2 приобретают суммарный положительный заряд.
\item При помещении стеклянных кубиков в электрическое поле наблюдается явление поляризации.
\item В электрическом поле кубики 3 и 4 приобретают суммарный отрицательный заряд.
\item После разделения кубик 1 имеет положительный заряд.
\item После разделения кубик 4 имеет положительный заряд.
\end{enumerate}
\begin{center}
\begin{tikzpicture}[book diagram,x=1cm,y=1cm]
\foreach \x/\y in {1.3/1.25,2.55/1.25,5.25/1.25,7.1/1.25}{
\fill[gray!8] (\x,\y) rectangle ++(1.25,1.25);
\foreach \h in {0.06,0.12,...,1.20}{\draw[gray!55,line width=0.2pt] (\x,{\y+\h})--++(1.25,0);}
\draw (\x,\y) rectangle ++(1.25,1.25);
}
\foreach \x/\y in {1.3/0,2.55/0,5.3/-0.75,7.1/-0.75}{
\fill[gray!30] (\x,\y) rectangle ++(1.25,1.25);
\foreach \h in {0.06,0.12,...,1.20}{\draw[gray!45,line width=0.2pt] (\x,{\y+\h})--++(1.25,0);}
\draw (\x,\y) rectangle ++(1.25,1.25);
}
\node at (1.925,1.875) {1};
\node at (3.175,1.875) {2};
\node at (1.925,0.625) {3};
\node at (3.175,0.625) {4};
\node at (5.875,1.875) {1};
\node at (7.725,1.875) {2};
\node at (5.925,-0.125) {3};
\node at (7.725,-0.125) {4};
\draw[-{book tip}] (-0.1,1.875)--(1.2,1.875);
\draw[-{book tip}] (-0.1,0.625)--(1.2,0.625);
\node[above] at (0.55,1.95) {$\vec E$};
\node[above] at (0.55,0.7) {$\vec E$};
\draw[dashed,gray] (4.48,-1.1)--(4.48,2.9);
\draw[dashed,gray] (4.72,-1.1)--(4.72,2.9);
\end{tikzpicture}
\end{center}
}{
\smallskip\noindent\textbf{Решение от GPT}\par
Стекло поляризуется, но суммарный заряд каждого стеклянного кубика остаётся нулевым: утверждение 2 верно, а 1 и 4 неверны. Общий заряд медных кубиков сохраняется: $q_3+q_4=0$, поэтому утверждение 3 неверно. Электроны переходят из кубика 4 в кубик 3 против направления поля; после разделения $q_3<0$, $q_4>0$, поэтому утверждение 5 верно. Ответ: 2, 5.
}

\end{document}
